\section{Task 1 - Classical Baseline}

The classical pipeline gets a test macro-F1 of 0.583 (0.613 on validation).

\subsection*{Method}
\begin{enumerate}
    \item \textbf{Language ID:} character 1--4-gram TF-IDF with logistic regression (L2, $C=1.0$), trained on the 100,446 tokens of the train split and their gold tags (\texttt{en}, \texttt{hi}, \texttt{univ}). Its predicted tags are used for every split.
    \item \textbf{Normalization:} from Dakshina, each Devanagari word's most frequent Roman spelling becomes the standard form, and every other spelling of that word maps to it. A spelling listed under several words follows the word it was seen with most often. That gives 21,130 variant $\rightarrow$ standard mappings. Only tokens tagged \texttt{hi} are changed (lookup in lower case), so English words stay as they are.
    \item \textbf{Sentiment:} word 1--3-gram + character 1--3-gram TF-IDF over the normalized tokens (182,648 features), with a Random Forest of 300 trees, max depth 20 and at least 5 samples per leaf.
\end{enumerate}

\subsection*{Results on the test split}
\begin{table}[H]
    \centering
    \begin{tabular}{@{}llc@{}}
        \toprule
        \textbf{Stage} & \textbf{Metric} & \textbf{Value} \\
        \midrule
        Language ID & Token accuracy & 0.871 \\
        Language ID & Macro-F1 (\texttt{en} / \texttt{hi} / \texttt{univ}) & 0.887 (0.78 / 0.88 / 0.999) \\
        Sentiment & Macro-F1 & 0.583 \\
        Sentiment & F1 (negative / neutral / positive) & 0.58 / 0.55 / 0.62 \\
        \bottomrule
    \end{tabular}
    \caption{Task 1 evaluation metrics on the Partition A test split.}
    \label{tab:task1_results}
\end{table}

\subsection*{Discussion}
\begin{itemize}
    \item The tagger almost never gets symbols and emojis wrong. Most of its mistakes are English vs. Hindi, on short words like ``to'', ``he'' or ``me'' that look the same in both languages.
    \item Normalization changed 6,122 of the 55,316 train tokens tagged as Hindi (11\%). Most changes merge real variants (\texttt{nahi} $\rightarrow$ \texttt{nahin}, \texttt{kyu} $\rightarrow$ \texttt{kyun}, \texttt{rha} $\rightarrow$ \texttt{raha}). Some are wrong, because the most frequent spelling isn't always the right word: \texttt{ji} $\rightarrow$ \texttt{ge} (from the acronym ``GE'') and \texttt{he} $\rightarrow$ \texttt{hey}.
    \item Neutral is the hardest class. In the confusion matrix, neutral tweets are confused with both other classes.
\end{itemize}
